The question this report answers is: \emph{what convergence rate does the
$\Linf$ error reach, when the closest point method is combined with an
angle rotation glue at a surface junction?}

The answer, on both surfaces, is \textbf{between first and second order,
and closer to first}. Table~\ref{tab:headline} collects it.

\begin{table}[htbp]
\centering
\caption{The headline numbers. Every rate is a least-squares slope of
  $\log\|e\|_\infty$ against $\log h$ over all the levels of that run.}
\label{tab:headline}
\small
\begin{tabular}{@{}llccl@{}}
\toprule
test & variant & levels & fitted rate & error at the finest $h$ \\
\midrule
\multicolumn{5}{@{}l}{\emph{Flipped cap on an ellipsoid}, $h=0.065
  \to 0.0241$} \\
\quad cut $z=0.390$, $\dk=3.07$ & exact $R$     & 6 & 1.26 & 6.44e-03 \\
\quad cut $z=0.455$, $\dk=2.68$ & exact $R$     & 6 & 1.80 & 5.62e-03 \\
\quad cut $z=0.520$, $\dk=2.35$ & exact $R$     & 6 & 1.47 & 3.75e-03 \\
\quad cut $z=0.585$, $\dk=2.07$ & exact $R$     & 6 & 1.23 & 3.24e-03 \\
\quad cut $z=0.390$, $\dk=3.07$ & \emph{no glue} & 6 & \textbf{0.10}
  & 2.84e-01 \\
\midrule
\multicolumn{5}{@{}l}{\emph{Half-twisted rectangular tube},
  $h=1/20 \to 1/44$} \\
\quad both corner curves & exact rotation & 5 & 1.43 & 4.42e-02 \\
\quad both corner curves & $d_{k2}$ rotation & 5 & 1.43 & 4.42e-02 \\
\bottomrule
\end{tabular}
\end{table}

\subsection*{The five findings}

\paragraph{1. The glue is what makes the solve converge.} On the flipped
ellipsoid, the same closest point method applied to the creased surface
\emph{without} the glue does not converge at all. Its fitted slope is
$0.10$ and its error stays near $0.3$ at every grid spacing. The glued
solve is 44 times more accurate at the same grid spacing, and it
converges. This is the largest effect in the whole study.

\paragraph{2. Neither surface reaches second order.} Every fitted rate
lies between $1.2$ and $1.8$. This is what the surface bound predicts
when the junction defect does not vanish: the term $h\,(D+J_2)$ is first
order and dominates the $h^2$ interior term. On the tube the
level-to-level rates fall steadily through the refinement
($2.23 \to 1.59 \to 0.86 \to 1.05$), so the fitted $1.43$ is an average
over a range in which the rate is still moving towards first order,
rather than an order the scheme holds.

\paragraph{3. The glue rotation does not have to be accurate.} The
rotations built only from closest point differences give the same
solution as the exact analytic rotation, everywhere in this study. On the
tube the two agree to $0.16$ per cent at worst, although the estimated
rotation is wrong by $1.40$~rad ($80^\circ$) on its worst row at the
coarsest level. On the ellipsoid the three variants agree to within 1 per
cent at the finest level of every cut. The glue has to be approximately
right; it does not have to be exact.

\paragraph{4. The transmission is not what limits the solve.} On both
surfaces the two branches agree with \emph{each other} across the
junction at a higher rate than either agrees with the exact solution:
$1.93$ and $1.71$ for the tube's two corner curves against $1.43$ for its
surface error, and $1.97$ for the ellipsoid's cross-branch gap against
$1.26$ for its surface error. The glue transmits consistently. The limit
is the geometry of the junction itself.

\paragraph{5. On the ellipsoid, the error is not a smooth function of
$h$.} Neighbouring grid spacings differ by a factor of three, and
level-to-level rates run from $-6.3$ to $+6.8$, while the linear solver
converges to $10^{-10}$ at every level. The cause is a signed
cancellation of the junction residual around the cut circle, whose
outcome depends on where that circle falls between grid planes. Only the
fitted slope is a rate. The tube does not show this, because its junction
curves are not plane circles and its glue angle varies along them.

\subsection*{What to be careful about}

\begin{itemize}
  \item \textbf{The fits are short.} The ellipsoid script's default is
        ten grid spacings, chosen to average over the scatter in
        finding~5. The runs here use six, because the finer levels of the
        default list cost more than the time available. The ellipsoid
        rates should be read as "first order, not second", not as
        three-digit numbers. See Section~\ref{sec:scatter}.
  \item \textbf{The $\dk=0$ control was not run.} The script can place
        the cut without flipping the cap, which makes the curvature
        mismatch vanish identically and should recover second order.
        That run would separate "the crease costs an order" from "the
        glue costs an order". It is the obvious next step.
  \item \textbf{The error level orders by $\dk$; the fitted rate does
        not.} The four cut heights give rates $1.26$, $1.80$, $1.47$ and
        $1.23$, which are not monotone in $\dk$. The errors at a fixed
        grid spacing are monotone. See Section~\ref{sec:dkappa}.
  \item \textbf{A bug blocks the script as committed.}
        \texttt{error\_surface\_figure} uses \texttt{thetaex}, which is
        never passed to it, so the ellipsoid script aborts after its
        first cut whenever \texttt{z\_3d} is not empty. The runs here use
        a copy with that argument added. See Section~\ref{sec:repro}.
\end{itemize}
